\paragraph{A precise limitation of the norm-ideal Massey formula.}
Sharifi's \emph{Massey products and ideal class groups},
J.~reine angew.~Math.~603 (2007), 1--33, has an author
version dated 26 July 2006 at
\url{https://www.math.ucla.edu/~sharifi/publ/massey.pdf}.
The page numbers in this paragraph refer to that forty-page
author version.  Section~4, pp.~17--18, fixes products with
consecutive copies of one Kummer class followed by a final
class.  The proper defining systems on pp.~21--22 require
a norm equation and an additional power condition.
Theorem~4.3 on p.~22 expresses the resulting value by
transgression, modulo the specified lower-product
indeterminacy.  Theorem~5.4 on p.~26, with its full proof
on pp.~26--27, gives the corresponding norm-ideal formula
for the prescribed ramification set.

These exact definitions, theorem statements and the complete
Theorem~5.4 proof were read for the present application.
KT-ORDINARY-CUBIC-SHUFFLE now proves the exact conversion
of the endpoint product $(u,v,u)$ to the consecutive repeat,
with explicit transport of its ordinary local primitives.
For $\mathbf F_p$-valued nonzero $u$, the binomial-block
normalization is explicit, but terminology alone does not
realize the specified last column by a norm representative.
The additional argument is KT-CUBIC-MATCHED-NORM-SYSTEM:
the actual class has zero absolute inflation, so cyclic
norm-residue exactness lifts its truncated cocycle through
the absolute regular module while preserving that last
column; Shapiro then supplies a matched norm representative.
KT-CUBIC-EXCISION retains the relative local correction in
the full rational class component, and
KT-CUBIC-P-ONLY-NORM applies the ideal formula with only
primes above $p$ inverted.  For a matched representative,
Sharifi's Proposition~4.2 and equation~(5.2) in the proof
of Theorem~5.4 give an exact ideal class, before the
lower-product quotient.  This repair uses the actual
zero absolute obstruction and does not assert that every
abstract binomial proper system is norm-realizable.
The norm-ideal formula still requires an arithmetic
evaluation.  A general finite-field-valued
$u$ also needs a scalar or cyclic reduction for this
particular norm formula.  No nonzero norm-ideal value,
injectivity assertion, or new arithmetic exclusion is
proved.  No reading of the entire forty-page paper or
reconstruction of its foundational proofs is claimed here.
